\section{Introduction}        
%%问题背景
\subsection{Problem Background}%%二级章节标题
Nowadays 
Two major problems are discussed in this paper, which are:
To solve this problem,we need to construct...
\begin{enumerate}[\quad\bf1.]  %%列出1、2、3点
    \item 
    \item 
\end{enumerate}
\begin{itemize}                %%以圆点方式给出各点
    \item First one...
    \item Second one ...
\end{itemize}
%%\subsubsection{Model 1}        %%三级章节标题
    \paragraph{}               %%一级段落。
    \subparagraph{}            %%二级段落。

%%文献综述,对现有文献在相似问题上的研究进行综述和评论
% \subsection{Literature Review}
% A literature\cite{2017-NCM-intro} say something about this problem ...%%引用

%%问题重述
\subsection{Clarification and Restatement}
In this problem we are given...,asked to...The data sets include...
We should only use these data to solve the following problems:
\begin{itemize}
    \item Based on the data given, the following problems need to be solved step by step first. 
    \begin{itemize}
        \item[-]Formulate an energy profile for each of the four states.
        \item[-]Design a model to show the evolving pattern of the energy profile. Find similarities and differences between the four states.
        \item[-]Based on the model, determine which state has the best energy usage profile in 2009.
        \item[-]According to the historical data and our conclusions, we predict the energy use of the states in 2025 and 2050, under the same condition.
    \end{itemize}
    \item Based on the above analysis, determine energy usage targets in 2025 and 2050 and make it as the goal of this new four states energy contracts, then provide at least three initiatives that will contribute to meet energy covenants. 
    \item Prepare a memo to the governor summarizing the energy situation for 2009, forecasting the future without policy changes and recommending targets.
\end{itemize}



%%工作介绍，简明叙述解题方法和主要成果
\subsection{Our work}
Our work follows the following workflow, as shown in the 
We do such things ...
We begin by establishing a model of ...
\begin{itemize}                %%以圆点方式给出各点
    \item We do...
    \item We do...
\end{itemize}